\documentclass[11pt]{article}
\usepackage{fontspec}
\usepackage{polyglossia}
\setmainlanguage{russian}
\setotherlanguage{english}
\setmainfont{Times New Roman}
\newfontfamily\cyrillicfont{Times New Roman}
\newfontfamily\cyrillicfontsf{Helvetica}
\newfontfamily\cyrillicfonttt{Menlo}
\usepackage{amsmath,amssymb}
\usepackage{booktabs}
\usepackage[margin=1in]{geometry}
\usepackage{enumitem}
\usepackage{hyperref}
\usepackage{xcolor}
\usepackage{tcolorbox}

\newtcolorbox{intuition}{colback=blue!5!white,colframe=blue!50!black,
  title=Интуиция,boxrule=0.5pt,fonttitle=\bfseries\small}
\newtcolorbox{tldr}{colback=yellow!10!white,colframe=orange!50!black,
  title=Главное,boxrule=0.5pt,fonttitle=\bfseries\small}

\title{Мощность теста Колмогорова--Смирнова в PhAIL:\\ от формулы из учебника к броуновскому мосту}
\author{}
\date{}

\begin{document}
\maketitle

\begin{abstract}
В аппендиксе статьи раньше стояла формула $N \approx 970$~--- столько эпизодов на руку требовалось бы для разрешения самой близкой пары моделей. Эмпирические кривые показывают, что хватает $N \approx 30$. Этот документ объясняет, откуда берётся это расхождение в 22~раза, и что мы поставили вместо старой формулы. Уровень: магистр CS без специализации в статистике.
\end{abstract}

\tableofcontents
\bigskip

%==============================================================
\section{В чём вопрос}
%==============================================================

В разделе \texttt{app:power} статьи мы приводим оценки нужного числа испытаний для разных тестов. Для теста Колмогорова--Смирнова (КС) старая формула давала $N \approx 970$ испытаний на руку~--- неприлично много на фоне нашего бюджета $\sim 30$ эпизодов на ячейку. При этом эмпирические кривые мощности (Рис.~\ref{fig:efficiency} в статье) показывают, что КС-тест \emph{реально} разрешает большинство близких пар при $N{=}25{-}30$.

Возникают три вопроса:
\begin{enumerate}[leftmargin=2em, itemsep=2pt]
\item Почему теоретическая формула давала $\sim 970$, а практика~--- $\sim 30$? Это формула неправильная или я её неправильно применил?
\item Если её заменить, на что? Можно ли получить теоретическую оценку, которая совпадает с эмпирикой?
\item Откуда конкретно берётся 22-кратный выигрыш? Это один большой эффект или сумма нескольких мелких?
\end{enumerate}

Этот документ отвечает на все три. По дороге придётся объяснить, что такое броуновский мост и почему он связан с КС-тестом, но я постараюсь делать это на пальцах.

%==============================================================
\section{Что такое тест Колмогорова--Смирнова}
%==============================================================

\subsection{Постановка}

Двухвыборочный тест КС спрашивает: "одинаковы ли два распределения?". У нас две выборки~--- $X_1, \ldots, X_n$ из распределения $F$ и $Y_1, \ldots, Y_m$ из распределения $G$. Нулевая гипотеза $H_0$: $F = G$. Альтернатива $H_1$: $F \neq G$ (в каком-то смысле).

\subsection{Статистика}

Нужно превратить "$F$ отличается от $G$" в одно число. КС берёт самое прямолинейное:

\[
D_{n,m} \;=\; \sup_t \bigl|F_n(t) - G_m(t)\bigr|,
\]

где $F_n, G_m$~--- эмпирические функции распределения (ЭФР). ЭФР строится тривиально: $F_n(t) = \frac{1}{n}\#\{i : X_i \leq t\}$. Это ступенчатая функция, поднимающаяся на $1/n$ при каждом $X_i$.

$D_{n,m}$~--- максимальное вертикальное расстояние между двумя ступеньками. Чем оно больше, тем сильнее различаются распределения.

\begin{intuition}
Нарисуйте две ступенчатые ЭФР на одной оси. КС-статистика~--- это самый длинный вертикальный отрезок, который можно провести между ними. Под $H_0$ ступеньки идут "ноздря в ноздрю", и это расстояние мало. Под $H_1$ они расходятся в каком-то месте, и расстояние большое.
\end{intuition}

\subsection{Распределение под $H_0$}

Это и есть магия КС: под $H_0$ распределение $D_{n,m}$ \emph{не зависит} от того, какое именно $F = G$ это было. Конкретно, при $n, m \to \infty$:

\[
\sqrt{\tfrac{nm}{n+m}}\, D_{n,m} \;\Rightarrow\; K,
\]

где $K$~--- так называемое \emph{распределение Колмогорова}. Его 95-я процентиль примерно равна $1.358$. То есть критерий КС на уровне 5\%: отвергнуть $H_0$, если $D_{n,m} > 1.358 \sqrt{(n+m)/(nm)}$.

Свойство "не зависит от $F$" называется distribution-free и означает, что одна и та же таблица критических значений работает для любого непрерывного $F$. Это не магия, а следствие того факта, что для непрерывного $F$ величина $F(X)$ распределена равномерно на $[0,1]$.

\subsection{Что КС делает в нашей статье}

В PhAIL мы сравниваем модели роботов по \emph{времени до успешного размещения предмета} $T$. У каждой модели $a$ есть распределение $F^{(a)}(t)$ "вероятность размещения за $\le t$ секунд". КС спрашивает: "разные ли распределения у моделей $a$ и $b$, или это одно и то же?".

Преимущество над сравнением одного числа (например, $F^{(a)}(30) - F^{(b)}(30)$): КС увидит различия в любой части кривой, а не только при $t = 30$. Если одна модель быстрее на быстрых задачах, а другая на медленных, $F(30)$ может это пропустить, а КС~--- нет.

%==============================================================
\section{Стандартная формула мощности и её предположения}
%==============================================================

\subsection{Что такое мощность}

Мощность теста~--- это $\Pr(\text{отвергнуть } H_0 \mid H_1)$. Если $H_1$ верна (распределения реально разные), мощность 80\%~--- это вероятность, что тест это поймает и выдаст $p < 0.05$. Чем больше выборка, тем больше мощность.

В планировании эксперимента вопрос ставится наоборот: \emph{какой нужен $N$, чтобы при заданном размере эффекта мощность была не ниже $1{-}\beta$?} (стандартное $\beta = 0.20$, то есть мощность 80\%).

\subsection{Формула, которая стояла в статье}

Для двухвыборочного КС с равными выборками размера $N$ стандартная асимптотическая формула:

\[
N \;\gtrsim\; \frac{2(c_\alpha + z_\beta)^2}{\Delta^2}
\;\approx\; \frac{9.68}{\Delta^2},
\qquad \alpha=0.05,\; \beta=0.20,
\]

где
\begin{itemize}[leftmargin=2em, itemsep=1pt]
\item $c_\alpha \approx 1.358$~--- критическое значение распределения Колмогорова на уровне $\alpha$,
\item $z_\beta \approx 0.84$~--- $(1-\beta)$-квантиль стандартного нормального распределения,
\item $\Delta = \sup_t |F^{(a)}(t) - F^{(b)}(t)|$~--- максимальное расстояние между истинными ФР.
\end{itemize}

При $\Delta = 0.10$ это даёт $N \approx 970$.

\subsection{Откуда берётся формула}

Вывод: под $H_1$ статистика $\sqrt{N/2}\, D_{N,N}$ ведёт себя примерно как нормальная величина с математическим ожиданием $\sqrt{N/2}\cdot\Delta$ и стандартным отклонением 1. Тест отвергает $H_0$, когда статистика превышает $c_\alpha$. Чтобы это происходило с вероятностью $1{-}\beta$:
\[
\sqrt{N/2}\cdot\Delta - c_\alpha \;\geq\; z_\beta
\;\;\Longleftrightarrow\;\;
N \;\geq\; \frac{2(c_\alpha + z_\beta)^2}{\Delta^2}.
\]

Запомните это рассуждение: оно опирается на два сильных приближения, и оба окажутся неточными в нашей ситуации.

%==============================================================
\section{Броуновский мост: что это и при чём он}
%==============================================================

Чтобы понять, где формула теряет точность, надо разобрать, что на самом деле происходит со статистикой $D_{N,N}$ под $H_1$.

\subsection{Винеровский процесс (краткое напоминание)}

\emph{Винеровский процесс} (или броуновское движение) $W(t)$~--- это случайный процесс на $[0, \infty)$ со свойствами:
\begin{itemize}[leftmargin=2em, itemsep=1pt]
\item $W(0) = 0$,
\item приращения независимы,
\item $W(t) - W(s) \sim \mathcal{N}(0, t-s)$ при $t > s$.
\end{itemize}
Картинка: непрерывная, очень "колючая" траектория, расходящаяся со временем (типичное отклонение $\sim \sqrt{t}$). Базовая модель случайного блуждания в непрерывном времени.

\subsection{Броуновский мост}

\emph{Броуновский мост} $B(t)$ на $[0, 1]$~--- это винеровский процесс, "пришпиленный" к нулю в обоих концах: $B(0) = B(1) = 0$. Формально:
\[
B(t) \;=\; W(t) - t\cdot W(1).
\]
То есть берём винеровский процесс и вычитаем линейный сдвиг, который заставляет его прийти в ноль при $t = 1$.

Свойства:
\begin{itemize}[leftmargin=2em, itemsep=1pt]
\item $B(t) \sim \mathcal{N}\bigl(0,\, t(1-t)\bigr)$ маргинально~--- максимум дисперсии $1/4$ в середине, нули по краям.
\item $\text{Cov}(B(s), B(t)) = \min(s, t) - st$.
\item $\sup_{t \in [0,1]} |B(t)|$ имеет распределение Колмогорова.
\end{itemize}

\begin{intuition}
Представьте маршрут случайного блуждания, которое \emph{обязано} вернуться в исходную точку за фиксированное время~--- например, цена акции, которая в конце дня должна закрыться там же, где открылась. Между концами она блуждает свободно, но с известным ограничением. Это и есть броуновский мост.
\end{intuition}

\subsection{Где КС встречается с мостом}

Ключевой факт: \emph{эмпирическая ФР, центрированная и масштабированная, ведёт себя как броуновский мост}. Точнее, для $X_1, \ldots, X_n$ из распределения $F$:
\[
\sqrt{n}\,\bigl(F_n(x) - F(x)\bigr) \;\Rightarrow\; B(F(x))
\quad \text{при } n \to \infty.
\]
(Это \emph{теорема Донскера}, и одна из главных теорем непараметрической статистики.)

Чем это полезно для КС: статистика $D_n = \sup_x |F_n(x) - F(x)|$ под $H_0$ ведёт себя как $\sup_x |B(F(x))|/\sqrt{n} = \sup_t |B(t)|/\sqrt{n}$~--- то самое супремум модуля моста, чьё распределение и есть распределение Колмогорова.

Для двух выборок: $\sqrt{n/2}\, D_{n,n} \Rightarrow \sup_t |B(t)|$ под $H_0$. Под $H_1$ (когда $F \neq G$):
\[
\sqrt{n/2}\, D_{n,n} \;\Rightarrow\; \sup_t \bigl|B(t) + \mu(t)\bigr|,
\quad \mu(t) = \sqrt{n/2}\,(F-G)(F^{-1}(t)).
\]
То есть к мосту добавляется \emph{сдвиг} $\mu(t)$, который и тащит статистику вверх.

\begin{tldr}
Под $H_1$ КС-статистика~--- это супремум модуля случайного процесса (моста), к которому добавлен \emph{детерминированный сдвиг}~--- кривая, описывающая, на сколько $F$ отличается от $G$ в каждой точке.
\end{tldr}

%==============================================================
\section{Где старая формула теряет точность: drift~--- это кривая, а не точка}
%==============================================================

Теперь к причине первой потери точности.

Стандартная формула $9.68/\Delta^2$ выводится в предположении, что $\sup_t |B(t) + \mu(t)|$ под $H_1$ имеет примерно гауссовское распределение со средним $\sqrt{n/2}\,\Delta$ (где $\Delta = \sup |F-G|$) и единичной дисперсией. Это всё равно что сказать: "статистика~--- это сдвиг на $\sqrt{n/2}\,\Delta$ в одной точке, плюс гауссовский шум".

\subsection{В чём проблема}

Реальный $\sup_t |B(t) + \mu(t)|$ зависит от \emph{всей кривой} $\mu(t)$, а не только от её максимума. Если $\mu$ "толстый" (то есть остаётся высоким на широком интервале $t$, или имеет несколько горбов), мост получает \emph{несколько шансов} превысить порог, а не один.

\begin{intuition}
Сравните:
\begin{itemize}[leftmargin=2em, itemsep=1pt]
\item $\mu_1$~--- горб шириной 0.01 высоты 0.12 (расхождение между ФР сидит в одной точке).
\item $\mu_2$~--- плато шириной 0.5 высоты 0.12 (распределения отличаются на половине области).
\end{itemize}
Максимум одинаковый, но во втором случае мост на этом плато колеблется случайно $\pm \sqrt{t(1-t)}$, и в \emph{какой-то} точке он почти наверняка зайдёт высоко~--- сумма сдвига и шума пробьёт порог. В первом случае такой шанс есть только в одной точке~--- сдвиг должен сам по себе быть достаточным.
\end{intuition}

В нашей задаче (время до успеха для разных моделей) $F-G$~--- именно широкая кривая: одна модель быстрее в окне $t \in [10, 30]$ секунд, потом догоняет, потом снова отстаёт и т.д. Старая формула этого не видит~--- она использует только максимальное значение $\Delta$.

\subsection{Цена ошибки: $\sim 2.5\times$}

Численно, для самой близкой пары (OpenPI vs.\ GR00T) с per-object $\Delta_o = 0.118$ на одном объекте:
\begin{itemize}[leftmargin=2em, itemsep=1pt]
\item Старая формула: $N \approx 9.68/0.118^2 \approx 695$ наблюдений на руку.
\item Симуляция моста с реальной кривой $F_o - G_o$: $N \approx 270$ наблюдений на руку даёт мощность 80\%.
\end{itemize}
Соотношение $\approx 2.5$. Это первая компонента итогового выигрыша.

%==============================================================
\section{Per-object vs pooled: почему мы считаем КС по-другому}
%==============================================================

Эмпирический тест в нашей статье делает не одно сравнение двух ФР, а \emph{четыре}~--- по одному на каждый из четырёх объектов (батарейки, ножницы, полотенца, деревянные ложки). Затем результаты усредняются: $\bar\Delta = \frac{1}{4}\sum_o \Delta_o$, где $\Delta_o = \sup_t |F^{(a)}_o(t) - F^{(b)}_o(t)|$.

Старая формула, наоборот, оперировала с $\Delta_\text{pool} = \sup_t |\bar F^{(a)}(t) - \bar F^{(b)}(t)|$, где $\bar F$~--- средняя ФР по объектам. Эти два подхода~--- сравнить-каждое-и-усреднить vs усреднить-и-сравнить~--- в общем случае дают разные числа.

\subsection{Почему усреднение ФР сжимает $\Delta$}

Расхождения по разным объектам происходят в \emph{разное время}: одна модель может быть быстрее на батарейках в $t = 10$ секунд, но медленнее на полотенцах в $t = 25$ секунд. Когда мы усредняем ФР, эти расхождения частично гасят друг друга~--- максимум $|\bar F^{(a)} - \bar F^{(b)}|$ оказывается заметно меньше, чем средний максимум по объектам.

Численно для наших трёх близких пар:

\medskip
\begin{tabular}{lcc}
\toprule
Пара & $\Delta_\text{pool}$ (sup на смеси) & $\bar\Delta$ (среднее по объектам) \\
\midrule
OpenPI vs GR00T & 0.035 & 0.122 \\
OpenPI vs ACT   & 0.118 & 0.142 \\
GR00T vs ACT    & 0.092 & 0.175 \\
\bottomrule
\end{tabular}
\medskip

Для самой близкой пары $\bar\Delta$ в \emph{3.5 раза} больше! Поскольку мощность зависит от $\Delta^2$, это уже само по себе кардинально меняет требуемое $N$.

\subsection{Кто из них правильный}

Оба корректны как тесты, но они отвечают на разные вопросы:
\begin{itemize}[leftmargin=2em, itemsep=1pt]
\item $\Delta_\text{pool}$: "если усреднить распределения по типам объектов, отличаются ли усреднённые распределения двух моделей?"
\item $\bar\Delta$: "в среднем по типам объектов, отличаются ли распределения двух моделей?"
\end{itemize}
Второй формулировкой пользуется наш бутстрэп~--- он считает per-object KS отдельно и усредняет. Старая формула применялась к первому, но "за неимением лучшего"~--- разумной альтернативы под рукой не было.

%==============================================================
\section{Macro-aggregation: ещё $3\times$ от усреднения четырёх тестов}
%==============================================================

Теперь второй большой выигрыш. Эмпирический тест~--- это не один тест, а \emph{среднее по $J = 4$ тестам}. Это даёт выигрыш по двум параметрам сразу: статистика становится менее шумной, а критическое значение опускается.

\subsection{Распределение средней статистики}

Каждая per-object $D_o$ под $H_0$ имеет среднее $\approx 0.87/\sqrt{n_o/2}$ и стандартное отклонение $\approx 0.26/\sqrt{n_o/2}$ (это моменты распределения Колмогорова). Если $D_o$ независимы между объектами (а они независимы, потому что эпизоды разных объектов~--- разные физические испытания), среднее $\bar D = \frac{1}{J}\sum D_o$ имеет:
\begin{itemize}[leftmargin=2em, itemsep=1pt]
\item то же среднее $\approx 0.87/\sqrt{n_o/2}$,
\item стандартное отклонение $\approx 0.26/(\sqrt{J}\sqrt{n_o/2})$~--- меньше в $\sqrt{J} = 2$ раза.
\end{itemize}
То есть критическое значение $\bar D$ при $\alpha = 0.05$ становится $\approx 1.08/\sqrt{n_o/2}$ (вместо $1.358/\sqrt{n_o/2}$ для одного теста)~--- понижение примерно на 25\%.

\subsection{Сигнал не падает}

Под $H_1$ среднее $\bar D \approx \bar\Delta = \frac{1}{J}\sum \Delta_o$. Если $\Delta_o$ примерно одинаковы по объектам (для нашей закрытой пары: 0.07, 0.12, 0.14, 0.15), $\bar\Delta$ остаётся в том же порядке, что и каждое отдельное $\Delta_o$. То есть \emph{сигнал такой же, шум меньше, порог ниже}~--- всё работает в нашу пользу.

\subsection{Итоговый выигрыш}

Численно: для самой близкой пары симуляция показывает, что для одного объекта (без усреднения) нужно $\approx 270$ эффективных наблюдений на руку для мощности 80\%, а для среднего по 4 объектам~--- $\approx 85$ наблюдений на руку \emph{на каждый объект}. Соотношение $\approx 3$.

\begin{tldr}
Усреднение KS по объектам вместо одного теста на смеси даёт $\sim 3\times$ выигрыш по нужному размеру выборки. Это не "магия данных"~--- это просто более чувствительная статистика, использующая структуру (несколько независимых объектов).
\end{tldr}

%==============================================================
\section{Кластеризация: эпизоды и операции}
%==============================================================

Осталось разобраться с тем, в каких единицах считать $N$. До сих пор мы говорили про "эффективные наблюдения на руку". Но эмпирический бюджет считается в \emph{эпизодах}~--- одна сессия работы робота. И эпизод~--- это не одно наблюдение.

\subsection{Структура данных}

Каждый эпизод $i$ содержит несколько успешных размещений; время между размещениями~--- наше $T$. В среднем по нашим данным $m \approx 4.4$ операций на эпизод. То есть если у нас $N_\text{cell}$ эпизодов на (модель, объект), то у нас примерно $N_\text{cell} \cdot m$ значений $T$.

\textbf{Можем ли мы просто умножить $N_\text{cell}$ на $m$ и получить эффективный размер выборки?} Нет: операции внутри одного эпизода \emph{сильно коррелированы}. Если модель в эпизоде "в ударе", все размещения быстрые; если "плывёт", все медленные. Каждая дополнительная операция в том же эпизоде даёт меньше информации, чем независимое наблюдение из нового эпизода.

\subsection{Эффект дизайна (design effect)}

Стандартная поправка из теории выборочных обследований: если у нас $n_\text{raw}$ наблюдений в кластерах среднего размера $m$ с внутрикластерной корреляцией $\rho$, то эффективный размер выборки для оценки среднего:
\[
n_\text{eff} \;=\; \frac{n_\text{raw}}{1 + (m-1)\rho},
\qquad D \equiv 1 + (m-1)\rho
\]
называется \emph{коэффициентом дизайна}.

Для наших данных: $m \approx 4.4$, $\rho \approx 0.7$ (мы измерили ICC на сырых $T$~--- внутрикластерная корреляция времён размещения внутри одного эпизода). Это даёт $D \approx 1 + 3.4 \cdot 0.7 \approx 3.4$. То есть $n_\text{eff} \approx 0.3 \cdot n_\text{raw}$.

\subsection{Почему $D$ для KS меньше}

Стандартная формула $D = 1 + (m-1)\rho$ выведена для \emph{среднего значения}. KS-тест работает не со средним, а с эмпирической ФР, то есть с индикаторами $\mathbf{1}\{T < t\}$. Корреляция этих индикаторов внутри кластера в общем случае \emph{меньше}, чем корреляция самих $T$.

\begin{intuition}
Представьте, что в эпизоде все четыре $T$~--- это одно и то же число $T_0$ (идеальная корреляция, $\rho = 1$). Тогда индикатор $\mathbf{1}\{T < 30\}$ для всех четырёх одинаков~--- кластер даёт ровно одну единицу информации про "выше или ниже 30". Вроде бы сильная корреляция совпадает.

Теперь представьте, что $T$ внутри эпизода скоррелированы, но не идеально~--- скажем, все четыре сидят в окне $\pm 5$ секунд от какого-то общего среднего. Корреляция самих $T$ высокая. Но индикаторы $\mathbf{1}\{T < t\}$ для $t$, далёких от этого окна, могут полностью совпадать (все равны 1 или все 0)~--- сильная корреляция; зато для $t$, близких к окну, разные операции могут оказаться по разные стороны порога~--- корреляция меньше. В среднем по $t$ корреляция индикаторов оказывается заметно меньше, чем корреляция $T$.
\end{intuition}

Поэтому стандартный $D = 3.4$ из ICC \emph{переоценивает} эффект кластеризации для KS. Чтобы найти правильный $D$, мы калибруем его эмпирически (см.\ след.\ раздел).

%==============================================================
\section{Симуляция броуновского моста как замена формуле}
%==============================================================

Чтобы получить правильную теоретическую оценку мощности, мы заменяем закрытую формулу на симуляцию. Алгоритм укладывается в $\sim 30$ строк кода.

\subsection{Генерация моста}

Нужно симулировать $B(t)$ на сетке $t \in [0,1]$. Стандартный приём: генерировать винеровский процесс $W$ инкрементами, потом вычесть линейный сдвиг.

\begin{enumerate}[leftmargin=2em, itemsep=2pt]
\item Зафиксировать сетку $0 = t_0 < t_1 < \ldots < t_K = 1$.
\item Генерировать инкременты $\Delta W_k \sim \mathcal{N}(0, t_k - t_{k-1})$ независимо.
\item Сложить кумулятивно: $W(t_k) = \sum_{j \leq k} \Delta W_j$.
\item Получить мост: $B(t_k) = W(t_k) - t_k \cdot W(t_K)$.
\end{enumerate}

В коде на numpy это $\sim 5$ строк (см.\ \texttt{build/explore\_ks\_power.py}).

\subsection{Алгоритм оценки мощности}

\begin{enumerate}[leftmargin=2em, itemsep=2pt]
\item Для каждого объекта $o$ оценить эмпирические ФР $F^{(a)}_o, F^{(b)}_o$ из всех имеющихся данных. Это наша лучшая догадка про "истинные" $F$ и $G$.
\item Для каждого объекта вычислить эффективное число операций на руку: $n_o = N_\text{cell} \cdot m_o / D$.
\item \textbf{Оценить распределение статистики под $H_0$:} симулировать $K$ мостов $B_o(t)$ для каждого объекта (без сдвига), вычислить $D_o^{\text{H}_0} = \sup_t |B_o(t)|/\sqrt{n_o/2}$, усреднить по объектам: $\bar D^{\text{H}_0} = \frac{1}{J}\sum_o D_o^{\text{H}_0}$. Получить 95-ю процентиль распределения $\bar D^{\text{H}_0}$~--- это критическое значение $q_{0.95}^{H_0}$.
\item \textbf{Оценить распределение статистики под $H_1$:} аналогично, но к мосту добавить детерминированный сдвиг $\mu_o(t) = \sqrt{n_o/2}\,(F^{(a)}_o - F^{(b)}_o)(t)$. Получить $D_o^{\text{H}_1} = \sup_t |B_o(t) + \mu_o(t)|/\sqrt{n_o/2}$, усреднить.
\item \textbf{Оценка мощности:} $\widehat{\text{Power}} = \Pr(\bar D^{\text{H}_1} > q_{0.95}^{H_0})$~--- доля симуляций, в которых статистика под $H_1$ превышает порог.
\end{enumerate}

\subsection{Калибровка $D$}

Остаётся параметр $D$~--- коэффициент дизайна для индикаторов. Мы калибруем его эмпирически: подставляем в симуляцию разные значения $D \in \{1.0, 1.5, \ldots, 5.0\}$, считаем предсказанную мощность для всех 6 точек $\{$3 близких пары$\} \times \{N_\text{cell} = 25, 30\}$, и сравниваем с реальной эмпирической мощностью (которую мы знаем из Рис.~3 в статье~--- она получена бутстрэпом).

Лучшее совпадение даёт $D \approx 2.25$. Это меньше, чем $D = 3.4$ из ICC на сырых $T$~--- именно из-за обсуждённого выше эффекта индикаторов.

\subsection{Валидация}

После калибровки на 6 точках мы предсказываем мощность для остальных $4 \cdot 3 = 12$ точек ($N_\text{cell} \in \{5, 10, 15, 20\}$, не использованных при калибровке) и сравниваем с эмпирикой:

\medskip
\begin{tabular}{lrrrr}
\toprule
Пара & $N_\text{cell}$ & Эмп. & Предск. & Ошибка \\
\midrule
OpenPI vs GR00T & 5  & 0.10 & 0.11 & +0.01 \\
OpenPI vs GR00T & 10 & 0.18 & 0.19 & +0.01 \\
OpenPI vs GR00T & 15 & 0.25 & 0.28 & +0.04 \\
OpenPI vs GR00T & 20 & 0.41 & 0.38 & --0.03 \\
OpenPI vs ACT   & 5  & 0.21 & 0.12 & --0.09 \\
OpenPI vs ACT   & 10 & 0.30 & 0.23 & --0.06 \\
OpenPI vs ACT   & 15 & 0.43 & 0.38 & --0.05 \\
OpenPI vs ACT   & 20 & 0.55 & 0.54 & --0.01 \\
GR00T vs ACT    & 5  & 0.20 & 0.20 & +0.01 \\
GR00T vs ACT    & 10 & 0.48 & 0.41 & --0.07 \\
GR00T vs ACT    & 15 & 0.68 & 0.63 & --0.05 \\
GR00T vs ACT    & 20 & 0.80 & 0.80 & 0.00 \\
\bottomrule
\end{tabular}
\medskip

Средняя абсолютная ошибка по всем 18 точкам (6 калибровочных + 12 валидационных)~--- около 3 процентных пунктов. Соотношение калибровочных к валидационным $1{:}2$, и ошибка не растёт на валидации~--- это разумная защита от перепиления.

%==============================================================
\section{Декомпозиция $14\times$: что откуда взялось}
%==============================================================

Возвращаемся к исходному вопросу: почему старая формула давала $N \approx 670$ на объект, а симуляция~--- $N_\text{cell} \approx 45$ на ячейку (что соответствует $\approx 92$ эффективным операциям на руку на объект). Соотношение в "родных" единицах симуляции (эффективные обс. на руку на объект): $670 / 92 \approx 7$. Если мерять в эпизодах на ячейку~--- ещё надо учесть $m/D$, итого $\sim 14\times$.

Откуда берётся фактор 7--14 (в зависимости от единиц измерения):

\begin{enumerate}[leftmargin=2em, itemsep=4pt]
\item \textbf{$2.5\times$ от моделирования сдвига как кривой.} Закрытая формула считает, что под $H_1$ статистика сдвигается на $\sqrt{n/2}\,\Delta$ "в одной точке". Реально сдвиг~--- это \emph{целая кривая} $\mu(t)$, и мост получает много шансов её перебить. Сравнение: на одном объекте симуляция требует $\sim 270$ наблюдений вместо формульных $\sim 670$. \emph{Этот выигрыш пропадает, если сдвиг действительно сосредоточен в одной точке (узкий горб).}

\item \textbf{$3\times$ от макро-усреднения по 4 объектам.} Тест на четырёх объектах с усреднением~--- это \emph{другая статистика}, не один тест на смеси. Усреднение $J$ независимых тестов даёт $\sqrt{J}$ снижение шума при том же сигнале~--- критическое значение опускается с $\approx 1.36$ до $\approx 1.08$ при $J=4$. Сравнение: с 4 объектами симуляция требует $\sim 85$ наблюдений на руку \emph{на объект} вместо $\sim 270$ для одного объекта. \emph{Этот выигрыш доступен только потому, что эмпирический тест уже устроен как macro-mean per-object KS~--- мы не "открыли" его, мы просто корректно его моделируем.}

\item \textbf{$2\times$ от пересчёта эффективных операций в эпизоды.} Это не настоящий выигрыш мощности, а просто смена единиц: $n_\text{eff} = N_\text{cell} \cdot m / D = N_\text{cell} \cdot 4.4 / 2.25 \approx 2 \cdot N_\text{cell}$. То есть один эпизод даёт примерно 2 эффективные независимые операции (а не 1 и не 4.4). Такой пересчёт превращает $\sim 92$ эффективных наблюдений в $\sim 45$ эпизодов на ячейку.
\end{enumerate}

\medskip
\begin{tabular}{lrr}
\toprule
Шаг & Эфф.~набл./рука/объект & Множитель \\
\midrule
{}[A] Закрытая формула $9.68/\bar\Delta^2$ & $\sim 670$ & $1.0\times$ \\
{}[B] Симуляция моста, один объект & $\sim 270$ & $2.5\times$ \\
{}[C] Симуляция, среднее по 4 объектам & $\sim 85$ & $\times 3 \to 7.5\times$ \\
{}[D] Перевод в эпизоды на ячейку & 45 эп.\,$\equiv$\,92 эфф.\,обс. & смена единиц \\
\bottomrule
\end{tabular}

\medskip
\begin{tldr}
Из 22-кратного расхождения между старой формулой ($N \approx 970$) и новой моделью ($N_\text{cell} \approx 45$):
\begin{itemize}[leftmargin=1em, itemsep=1pt]
\item часть~--- бухгалтерская (правильно ли считать в эпизодах или в операциях),
\item $\sim 3\times$~--- от усреднения per-object KS вместо КС на смеси (это сильнее как тест),
\item $\sim 2.5\times$~--- от честного моделирования $H_1$-распределения статистики (старая формула слишком грубая).
\end{itemize}
Оба настоящих выигрыша~--- это не "новые трюки", а корректное описание того, что эмпирический тест и так делает. Новая модель просто даёт ему правильную теоретическую интерпретацию.
\end{tldr}

%==============================================================
\section{Итоги: что мы предсказываем и почему этому стоит верить}
%==============================================================

Подставив калибровочные значения $D = 2.25$, $m = 4.4$ и реальные $F_o, G_o$ для каждой из трёх близких пар, симуляция предсказывает $N_{0.8}$~--- минимальное число эпизодов на ячейку для достижения мощности 80\%:

\medskip
\begin{tabular}{lrr}
\toprule
Пара & $N_{0.8}$ предсказание & Эмпирика при $N_\text{cell}=30$ \\
\midrule
GR00T vs.\ ACT     & 25 эп./яч. & 0.95 (уже выше 0.8) \\
OpenPI vs.\ ACT    & 30 эп./яч. & 0.78 (на пороге) \\
OpenPI vs.\ GR00T  & 45 эп./яч. & 0.62 (не достигнуто) \\
\bottomrule
\end{tabular}
\medskip

Для двух пар, где эмпирика уже разрешает~--- предсказание совпадает (модель не противоречит наблюдаемому). Для самой близкой пары, которая в текущем бюджете не разрешается~--- предсказание говорит, что увеличение бюджета на $\sim 50\%$ (с 30 до 45 эпизодов на ячейку) её разрешит. Это \emph{не} порядок величины, как намекала старая формула.

Почему этому стоит верить:
\begin{itemize}[leftmargin=2em, itemsep=2pt]
\item Модель внутренне согласована: использует ту же статистику, что и эмпирический тест (per-object macro KS), и тот же бутстрэп для $H_0$.
\item Калибрация одного параметра ($D = 2.25$) на 6 точках, валидация на 12~--- не переобучаемся.
\item Декомпозиция $14\times$ на физически интерпретируемые компоненты ($2.5\times$ от drift, $3\times$ от macro)~--- никаких "магических" коэффициентов.
\end{itemize}

%==============================================================
\section{Что отвечать на типичные вопросы}
%==============================================================

\paragraph{"Этот $D = 2.25$~--- подгонка?"}
Один параметр калибруется на 6 точках, проверяется на 12~--- это нормальное соотношение. Альтернатива (использовать $D$ из ICC, $\approx 3.4$) \emph{переоценивает} эффект кластеризации, потому что ICC меряет корреляцию $T$, а KS работает с индикаторами $\mathbf{1}\{T < t\}$, у которых корреляция меньше. У нас просто нет красивой замкнутой формулы для $D$ под индикаторы~--- поэтому пинаем эмпирически.

\paragraph{"Зачем теоретическая модель, если есть эмпирические кривые?"}
Эмпирические кривые мощности доступны для $N \in [5, 30]$~--- это то, на чём бутстрэп подсчитан. Для самой близкой пары мощность не достигла 0.8 в этом диапазоне. Без модели мы не знаем, нужно ли 35, 50 или 200 эпизодов~--- модель даёт обоснованное "около 45".

\paragraph{"Зачем броуновский мост, если всё равно симулируем?"}
Бутстрэп~--- это золотой стандарт для $N \in [5, 30]$, потому что он использует реальные данные. Симуляция моста~--- это \emph{параметрическая модель}, которая позволяет экстраполировать наружу. Кроме того, моделирование через мост даёт асимптотическую интерпретацию: понятно, какой именно стохастический процесс описывает статистику, и как он зависит от $F$ и $G$.

\paragraph{"А ваше непрерывное приближение моста~--- ОК для дискретных $T$?"}
Под $H_0$ ничего не страдает: и наблюдения, и критическое значение получаются из одной и той же симуляции с одним и тем же приближением, поэтому смещения сокращаются. Эмпирические 3 п.п.\ совпадение на всех 18 точках~--- лучшее доказательство, что приближение нам не вредит.

\paragraph{"Почему именно per-object усреднение, а не stratified KS какой-нибудь?"}
Это выбор эмпирической процедуры~--- мы моделируем то, что есть. Stratified KS со взвешиванием по числу испытаний может быть ещё мощнее, но это другая методология. Сейчас мы строим теорию для уже существующего теста.

\bigskip
\begin{center}\rule{0.4\textwidth}{0.4pt}\end{center}

\paragraph{Полные числа и код:} \texttt{build/explore\_ks\_power.py},  \texttt{build/explore\_ks\_power\_v3.py}; выходные файлы~--- в \texttt{out/explore/ks\_power/}.

\end{document}
